

\begin{frame}
\frametitle{Reproducibility evaluation levels: guideline for authors}
\textbf{Availability of the code \\}
\begin{tabular}{|p{10cm}|p{1.5cm}|}
%\begin{tabular}{|c|c|}
  \hline
  \cellcolor{bl1}  \cellcolor{bl1} Description & \cellcolor{bl1} Stamp \\
\hline
code available on a public forge + license, authors, contributors files  & \textcolor{green1}{\ding{78}} \\
commit number and branche corresponding to the publication & \\
\hline
+ code available on Zenodo & \textcolor{green1}{\ding{78}\ding{78}} \\
\hline
+ Software Heritage Identifier (SWHID) &  \textcolor{green1}{\ding{78}\ding{78}\ding{78}}\\

\hline
+ HAL deposit & \textcolor{green1}{\ding{78}\ding{78}\ding{78}\ding{78}} \\ 
\hline
\end{tabular}
~\\

\textbf{Reproducibility \\}
\begin{tabular}{|p{10cm}|p{1cm}|p{1.7cm}|}
\hline
\cellcolor{bl1}  \cellcolor{bl1} Description & \cellcolor{bl1} Stamp & \cellcolor{bl1} Requirement \\
\hline
re-runable code, description of the workflow and test case(s) &  \textcolor{yell1}{\ding{115}} & \textcolor{green1}{\ding{78}}\\ 
\hline
+ workflow manager for automatic execution  &  \textcolor{yell1}{\ding{115}\ding{115}} & \textcolor{green1}{\ding{78}}\\ 
\hline
+ reproducible software environment & \textcolor{yell1}{\ding{115}\ding{115}\ding{115}} & \textcolor{green1}{\ding{78}} \\
\hline
+ regression, test cases and documentation in CI & \textcolor{yell1}{\ding{115}\ding{115}\ding{115}\ding{115}} &  \textcolor{green1}{\ding{78}} \\
 \hline
\end{tabular}
\end{frame}

%%%%%%%%%%%%%%%%%%%%%%%%%%%%%%%%%%%%%%%%%%

\begin{frame}
\frametitle{Reproducibility evaluation levels: guideline for authors}
\textbf{Re-usability for further projects}
\begin{tabular}{|p{10cm}|p{1cm}|p{1.7cm}|}
\hline
\cellcolor{bl1}  \cellcolor{bl1} Description & \cellcolor{bl1} Stamp & \cellcolor{bl1} Requirement \\
\hline
All functions documentation inside the code & \textcolor{purple1}{\ding{112}} &  \textcolor{yell1}{\ding{115}} \\
Description of each element in a Readme + workflow description for execution &  \textcolor{purple1}{\ding{112}\ding{112}} &\\
User documentation with examples &  \textcolor{purple1}{\ding{112}\ding{112}\ding{112}}&\\
Developper documentation (``How to contribute ?'' guide) &  \textcolor{purple1}{\ding{112}\ding{112}\ding{112}\ding{112}}&\\
\hline
\end{tabular}

\end{frame}


\begin{frame}
  Resources (summary)
\end{frame}



\begin{frame}
\frametitle{ACM badges 1/2}

One badge by element. Can be attributed independently. Can be cumulated.\\
"artifact"= digital object created by the authors to be used as part of the study or generated by the experiment itself.

%\begin{tabular*}{0.5\linewidth}{@{\extracolsep{\stretch{1}}}lp{8cm}@{}}
\begin{tabular}{|p{5cm}|p{9cm}|}
  \hline
 & \cellcolor{bl1}Validated elements \\
  \hline 
  Badge  \textbf{artifacts evaluated}: functional \\ \textit{available for the reviewer}  & The artifacts associated with the research are found to be documented, consistent, complete, exercisable, and include appropriate evidence of verification and validation \\
\hline
  Badge \textbf{artifacts evaluated}: re-usable \\ \textit{available for the reviewer}   &  The artifacts associated with the paper are of a quality that significantly exceeds minimal functionality. That is, they have all the qualities of the Artifacts Evaluated – Functional level, but, in addition, they are very carefully documented and well-structured to the extent that reuse and repurposing is facilitated. In particular, norms and standards of the research community for artifacts of this type are strictly adhered to. \\
  \hline
\end{tabular}
\end{frame}





\begin{frame}
\frametitle{ACM badges 2/2}

\begin{tabular}{|p{5cm}|p{9cm}|}
   \hline 
Badge \textbf{artifacts available} \\ \textit{publically} & Author-created artifacts relevant to this paper have been placed on a publically accessible archival repository. A DOI or link to this repository along with a unique identifier for the object is provided. \\
  \hline
Badge \textbf{Results reproduced}& The main results of the paper have been obtained in a subsequent study by a person or team other than the authors, using, in part, artifacts provided by the author. \\
  \hline
Badge \textbf{Results replicated} & The main results of the paper have been independently obtained in a subsequent study by a person or team other than the authors, without the use of author-supplied artifacts. \\ 
  \hline 
\end{tabular}
https://www.replicabilitystamp.org/: Papers with the Replicability Stamp are cited 15$\%$ more than average one year after publication, and 33$\%$ more two years after publication.

[Calculated based on Clarivate Web Of Science (WoS) citation counts in 2020]

\end{frame}


\begin{frame}
\frametitle{Benureau / Rougier}
\setcellgapes{0.3em}
\vspace{-0.2cm}
\begin{tabular}{|p{1cm}|p{2cm}|p{10cm}|}
  \hline
\cellcolor{bl1}Level & \cellcolor{bl1} State & Definition \cellcolor{bl1}\\
  \hline
$R^1$   & \textbf{re-runnable} & A re-runnable code is one that can be run again when needed, and in particular more than the one time that was needed to produce the results. \\
  \hline
 $R^2$   & \textbf{repeatable} & The code is running and producing the expected results. The next step is to make sure that you can produce the same output over successive runs of your program. In other words, the next step is to make your program deterministic, producing repeatable output. \\
  \hline
  $R^3$   & \textbf{reproducible} & A result is said to be reproducible if another researcher can take the original code and input data, execute it, and re-obtain the same result. \\
  \hline
$R^4$    & \textbf{re-usable} & Making your program reusable means it can be easily used, and modified, by you and other people, inside and outside your lab (comments, documentation,workflow)  \\ 
  \hline
 $R^5$    & \textbf{replicable} & Replicability is the implicit assumption that an article that does not provide the source code makes: that the description it provides of the algorithms is sufficiently precise and complete to re-obtain the results it presents. \\ 
  \hline 
\end{tabular}
\end{frame}




%-------------------------------------------------------------
\begin{frame}
\frametitle{CODE beyond FAIR 1/2}

Recommandations : progressive ROADMAP

\begin{tabular}{|l|l|}
  \hline
 \cellcolor{bl1} OPEN & \cellcolor{bl1}\\
    \hline 
• Publish your source code on a public forge & mandatory \\
• Save your repository on dedicated archive & mandatory \\
• License your code with \textcolor{red}{an open license} & strongly recommended \\
• Declare authorship and rightholders & recommended \\
  \hline 
 \cellcolor{bl1} DOCUMENT & \cellcolor{bl1}\\
    \hline 
• Choose meaningful names & recommended \\
• Comment code & recommended \\
• Provide examples, notebooks or tutorials & recommended \\
• Describe API & optional \\
  \hline 

\end{tabular}
\end{frame}






\begin{frame}
\frametitle{CODE beyond FAIR 2/2}

Recommandations : progressive ROADMAP

\begin{tabular}{|l|l|}
  \hline
\cellcolor{bl1} EXECUTE & \cellcolor{bl1}\\
    \hline 
• List software and hardware dependencies & recommended \\
• Provide a computational environment & optional \\
• Implement a test suite & optional \\
• Show real-life usage example with expected results & optional \\
  \hline 
\cellcolor{bl1}  COLLABORATE & \cellcolor{bl1}\\
    \hline 
• Respond to issues and feature requests & optional \\
• Describe maintenance, features and support limits & optional \\
• Explain how to contribute & optional \\
• Build and animate a community &  optional \\
  \hline 
\end{tabular}
\end{frame}


\begin{frame}[shrink=5]
  \frametitle{Other ressources}

  \begin{itemize}
  \item https://www.replicabilitystamp.org/
  \item \href{https://doc.hal.science/ressources/essentiels/CCSD_essentiels_logiciels.fr.pdf}{HAL - Le dépôt de logiciels}
  \item  https://inria.hal.science/hal-04697010/file/final-author-version.pdf
  \item  https://inria.hal.science/hal-04930405v1/document
  \item  https://computo-journal.org/review/
  \item  https://www.frontiersin.org/journals/neuroinformatics/articles/10.3389/fninf.2017.00069/full
  \end{itemize}

\end{frame}

%%% Local Variables:
%%% mode: latex
%%% TeX-master: "main"
%%% End:
